% TikZ diagrams for the working notes.  Each is a macro so it can be placed in a figure.

% ---------------------------------------------------------------- Penrose diagram
% Evaporating black hole formed by a collapsing null shell (null lines at 45 degrees),
% zoomed on the late-time region.
\newcommand{\PenroseSlice}{%
\begin{tikzpicture}[scale=1.15, >=Stealth, font=\small,
    quantum/.style={decorate, decoration={snake, amplitude=1.1pt, segment length=5pt, post length=3pt}, ->}]
  \clip (-1.0,2.75) rectangle (6.9,9.45);
  \coordinate (im) at (0,0);   \coordinate (iz) at (6,6);   \coordinate (ip) at (3,9);
  \coordinate (E) at (3,7.5);  \coordinate (S0) at (0,7.5);
  \fill[black!4] (im) -- (iz) -- (ip) -- (E) -- (S0) -- cycle;
  \fill[black!11] (0,4.5) -- (E) -- (S0) -- cycle;                       % black hole region
  \draw[thick] (im) -- (S0) node[pos=0.5, left] {$r=0$};
  \draw[thick] (E) -- (ip);
  \draw[thick] (im) -- (iz) node[pos=0.8, below right] {$\mathscr I^-$};
  \draw[thick] (iz) -- (ip) node[pos=0.4, above right] {$\mathscr I^+$};
  \draw[thick, decorate, decoration={zigzag, segment length=5pt, amplitude=1.5pt}] (S0) -- (E)
     node[midway, above=3pt] {singularity};
  \draw[very thick, black!50] (3,3) -- (0,6) node[pos=0.3, left=3pt, font=\scriptsize, black!60] {collapsing shell};
  \draw[thick, dashed] (0,4.5) -- (E);
  \node[font=\scriptsize, rotate=45] at (0.95,4.95) {event horizon};
  % one Hawking pair: outgoing quantum to I+, partner behind the horizon
  \draw[quantum, black!75] (2.55,6.9) -- (3.65,8.0);
  \draw[quantum, black!45] (2.3,7.0) -- (1.9,7.4);
  % Cauchy slice
  \coordinate (Ip) at (5.1,6.9); \coordinate (B) at (1.8,6.3); \coordinate (F) at (0.25,5.75);
  \draw[line width=1.7pt, blue!70!black] (Ip) -- (B);
  \draw[line width=1.7pt, red!70!black] (B) -- (F);
  \fill (B) circle (2.4pt); \fill[black!65] (F) circle (2.1pt);
  \node[blue!70!black, below] at (3.6,6.55) {$r\;/\;R$};
  \node[red!70!black] at (0.85,6.35) {$\ell$};
  \draw[black!60] (F) -- (-0.35,5.95) node[left, black] {$f$};
  \node at (2.0,6.0) {$B$};
  \node[blue!70!black, font=\scriptsize] at (5.55,7.25) {slice $\Sigma_t$};
  \node[right] at (iz) {$i^0$}; \node[above] at (ip) {$i^+$};
\end{tikzpicture}}

% ---------------------------------------------------------- two compressions
\newcommand{\CompressionPanel}[8]{% xshift, title, big, set, map, bigdim, small, smalldim
    \begin{scope}[xshift=#1 cm]
      \node[font=\small\bfseries] at (1.6,3.85) {#2};
      \draw[thick, fill=black!3] (1.6,1.75) ellipse (1.55 and 1.35);
      \node at (0.75,2.45) {$#3$};
      \node[font=\scriptsize, black!60] at (1.6,3.38) {$\dim=#6$};
      \fill[orange!35, draw=orange!70!black] plot[smooth cycle, tension=0.8]
        coordinates {(1.0,1.6) (1.5,2.15) (2.15,2.0) (2.3,1.45) (1.7,1.2)};
      \node at (1.65,1.65) {$#4$};
      \fill[pattern=north east lines, pattern color=black!35] (1.0,0.75) ellipse (0.55 and 0.25);
      \node[font=\scriptsize, black!60, anchor=west] at (1.6,0.55) {null states};
      \draw[thick, fill=black!3] (1.6,-1.25) ellipse (0.8 and 0.5);
      \node at (1.6,-1.25) {$#7$};
      \node[font=\scriptsize, black!60, anchor=west] at (2.45,-1.25) {$\dim=#8$};
      \draw[->, very thick] (1.6,0.35) -- node[right] {$#5$} (1.6,-0.7);
    \end{scope}}
\newcommand{\TwoCompressions}{%
\begin{tikzpicture}[>=Stealth, font=\small]
  \CompressionPanel{0}{black hole interior}{\mathcal H_\ell}{\Sigma_C}{V}{e^{S_\ell}}{\mathcal H_B}{e^{S_{\rm BH}}}
  \CompressionPanel{6.4}{network layer}{\mathbb R^m}{\Sigma_k}{W}{m}{\mathbb R^n}{n}
  \node[align=center, font=\small] at (4.8,-2.3)
    {$\sup_{x,y\in\Sigma}\big|\langle Vx,Vy\rangle-\langle x,y\rangle\big| \;\lesssim\; \sqrt{\log|\Sigma|\,/\,\beta\,\dim}$\\[1pt]
     \scriptsize one theorem: $\beta=2$ and $\Sigma=$ low complexity (left); $\beta=1$ and $\Sigma=$ sparse (right)};
\end{tikzpicture}}

% ------------------------------------------------- toy model of superposition
\newcommand{\TMSPicture}{%
\begin{tikzpicture}[>=Stealth, font=\small, scale=1.0]
  % left: bipartite features -> neurons
  \foreach \i in {1,...,8} \node[circle, draw, minimum size=11pt, inner sep=0pt, fill=white] (f\i) at (0.55*\i,2.2) {};
  \foreach \i in {2,7} \node[circle, draw, minimum size=11pt, inner sep=0pt, fill=orange!45] at (f\i) {};
  \foreach \j in {1,2,3} \node[rectangle, draw, minimum size=12pt, inner sep=0pt, fill=black!6] (n\j) at (0.55+1.1*\j,0.6) {};
  \foreach \i/\j/\w in {1/1/0.8, 2/1/0.4, 2/2/1.2, 3/2/0.6, 4/3/0.9, 5/1/0.5, 5/3/0.5, 6/2/0.8, 7/3/1.2, 7/1/0.4, 8/2/0.6, 3/3/0.3, 6/1/0.3, 8/3/0.5, 1/2/0.2, 4/1/0.3}
     \draw[line width=\w pt, black!60] (f\i) -- (n\j);
  \node[anchor=east, font=\scriptsize] at (0.3,2.2) {features $x\in\mathbb R^m$};
  \node[anchor=east, font=\scriptsize] at (1.4,0.6) {neurons $h=Wx$};
  \node[font=\scriptsize, black!60, align=center] at (2.5,-0.05) {each neuron reads many features\\(polysemantic)};
  % right: pentagon of feature directions in R^2
  \begin{scope}[xshift=6.6cm, yshift=1.3cm]
    \draw[black!30] (-1.3,0) -- (1.3,0); \draw[black!30] (0,-1.3) -- (0,1.3);
    \foreach \k in {0,2,3,4} \draw[->, thick] (0,0) -- ({90+72*\k}:1.15);
    \draw[->, thick, orange!80!black] (0,0) -- (162:1.15);
    \node[font=\scriptsize] at (0,-1.6) {$m=5$ features in $n=2$ neurons};
    \node[font=\scriptsize, orange!80!black] at (-1.45,0.55) {$w_i$};
  \end{scope}
\end{tikzpicture}}

% ----------------------------------------------------- sender/receiver/monitor
\newcommand{\ChannelPicture}{%
\begin{tikzpicture}[>=Stealth, font=\small, node distance=1.25cm,
    box/.style={draw, thick, rounded corners=2pt, minimum height=0.9cm, align=center, fill=black!3}]
  \node[box] (x) {features\\$x\in\{0,1\}^d$};
  \node[box, right=0.9cm of x] (s) {sender\\$h=\tanh(Ax+a)$};
  \node[box, right=0.9cm of s, fill=blue!6] (h) {message $h\in\mathbb R^n$\\[-1pt]\scriptsize public frame $=$ coordinates};
  \node[box, right=0.9cm of h] (r) {receiver\\MLP};
  \node[box, right=0.9cm of r] (xh) {$\hat x$};
  \draw[->, thick] (x) -- (s);
  \draw[->, thick] (s) -- (h);
  \draw[->, thick] (h) -- node[above, font=\scriptsize] {$+\sigma\xi$} (r);
  \draw[->, thick] (r) -- (xh);
  \node[box, below=1.0cm of h, fill=orange!12, text width=5.6cm] (m)
     {linear monitor $\mathrm{sigmoid}(Mh+m_0)$\\[-1pt]\scriptsize frozen $\;\to\;$ recalibrated (scale, bias) $\;\to\;$ refit};
  \draw[->, thick, orange!70!black] (h) -- (m);
  \draw[->, dashed, thick, orange!70!black] (m.west) -| (s.south);
  \node[font=\scriptsize, orange!70!black, anchor=north] at ($(s.south)+(0,-1.55)$) {optional: monitor in the loss ($\mu$)};
\end{tikzpicture}}

% ---------------------------------------------- legibility levels (ladder)
\newcommand{\MonitorLadder}{%
\begin{tikzpicture}[font=\small, >=Stealth]
  \foreach \y/\lab/\desc in {3/receiver/{what the channel carries (task F1)},
                             2/refit/{best linear public decoder, retrained},
                             1/recalibrated/{frozen directions; new scale and bias},
                             0/frozen/{decoder fixed at the legible checkpoint}} {
     \draw[thick] (0,0.95*\y) -- (1.4,0.95*\y);
     \node[anchor=east] at (-0.1,0.95*\y) {\lab};
     \node[anchor=west, font=\scriptsize, black!65] at (1.55,0.95*\y) {\desc};
  }
  \draw[<->, thick, blue!60!black] (6.9,0.05) -- node[right, font=\scriptsize, align=left] {calibration drift} (6.9,0.9);
  \draw[<->, thick, red!60!black] (6.9,1.0) -- node[right, font=\scriptsize, align=left] {frame (direction) drift} (6.9,1.85);
  \draw[<->, thick, black!70] (6.9,1.95) -- node[right, font=\scriptsize, align=left] {hidden channel} (6.9,2.8);
\end{tikzpicture}}
